\documentclass[12pt]{article}

\usepackage[utf8]{inputenc}
\usepackage[T1]{fontenc}
\usepackage[brazil]{babel}

\usepackage{geometry}
\geometry{margin=2.5cm}
\usepackage{setspace}
\onehalfspacing

\usepackage{amsmath, amssymb, amsthm}
\usepackage{mathtools}

\numberwithin{equation}{section}

\theoremstyle{plain}
\newtheorem{theorem}{Teorema}[section]
\newtheorem{proposition}[theorem]{Proposição}
\newtheorem{lemma}[theorem]{Lema}
\newtheorem{corollary}[theorem]{Corolário}

\theoremstyle{definition}
\newtheorem{definition}[theorem]{Definição}
\newtheorem{example}[theorem]{Exemplo}
\newtheorem{exercise}{Exercício}[section]

\theoremstyle{remark}
\newtheorem*{remark}{Observação}

\usepackage{hyperref}
\usepackage{csquotes}

\usepackage[
    backend=biber,
    style=authoryear,
    maxcitenames=2
]{biblatex}

\addbibresource{/mnt/c/Users/nicho/Documents/nicholasvoltani.github.io/bibliography.bib}

\newcommand{\R}{\mathbb{R}}
\newcommand{\pref}{\succeq}
\newcommand{\strictpref}{\succ}
\newcommand{\indiff}{\sim}

% ------------------------------------------------
% Title (fill manually)
% ------------------------------------------------
\title{}
\author{}
\date{}

\begin{document}

\maketitle

\hypertarget{exercuxedcio-2-utilidade-ces}{%
\section{Exercício 2 --- Utilidade
CES}\label{exercuxedcio-2-utilidade-ces}}

!500

Tem-se que a Função de Utilidade CES \[
u(x) = (\alpha_{1} x_{1}^\rho + \alpha_{2}x_{2}^\rho)^{1/\rho}
\] é linear quando \(\rho = 1\), e converge para a Função de
Cobb-Douglas quando \(\rho \to 0\).
\textbf{Suponha-se também que \sum_{i} \alpha_{i} = 1.}

Para verificar isso, podemos fazer a sequência de operações{[}\^{}6{]}
\[
\ln \to \lim_{ \rho \to 0 } \to \exp \iff \lim_{ \rho \to 0 } 
\]

\[
\tilde{u}_{\rho}(x) \equiv \ln(u(x)) = \frac{1}{\rho} \ln(\alpha_{1} x_{1}^\rho + \alpha_{2}x_{2}^\rho)
\] Tomando o limite, teríamos uma situação \(\frac{0}{0}\), i.e.~regra
de L'Hôpital: \[
\lim_{ \rho \to 0 } \tilde{u}_{\rho}(x) = \frac{\frac{ \partial  }{ \partial \rho } \ln(\alpha_{1} x_{1}^\rho + \alpha_{2}x_{2}^\rho)}{1}
\]

Porém, \textbf{a derivada não é em relação a x}! \[
\frac{ \partial x^\rho }{ \partial \rho } = \frac{ \partial  }{ \partial \rho } \exp(\ln x^\rho) = \frac{ \partial  }{ \partial \rho } \exp(\rho \ln x) = \ln(x) \exp(\ln x^\rho) = \ln(x) x^\rho
\]

Portanto, o limite fica \[
\begin{align*}
\lim_{ \rho \to 0 } \tilde{u}_{\rho}(x) &= \lim_{ \rho \to 0 } \frac{\alpha_{1} \ln(x_{1}) x_{1}^\rho + \alpha_{2} \ln(x_{2}) x_{2}^\rho}{\alpha_{1} x_{1}^\rho + \alpha_{2} x_{2}^\rho} \\
&= \alpha_{1} \ln x_{1} + \alpha_{2} \ln x_{2}
\end{align*}
\] Re-exponenciando, temos que \[
\lim_{ \rho \to 0 } u(x) = x_{1}^{\alpha_{1}} x_{2}^{\alpha_{2}}
\]

Uma visualização dessa convergência: !500 Fonte:
\href{https://github.com/nicholasvoltani/Manimations/tree/master/ConstantElasticitySubstitution_Function}{Eu
mesmo!}

\hypertarget{exercuxedcio-3}{%
\section{Exercício 3}\label{exercuxedcio-3}}

!500

\hypertarget{demandas-walrasianas-da-ces}{%
\subsubsection{Demandas walrasianas da
CES}\label{demandas-walrasianas-da-ces}}

Pode-se resolver a Demanda Walrasiana para a Função de Utilidade CES em
geral. Para isso, compensa resolver o problema de otimização para a nova
utilidade{[}\^{}4{]} \[
\tilde{u}(x) \equiv \rho \, u(x)^\rho = \rho (\alpha_{1} x_{1}^\rho + \alpha_{2} x_{2}^\rho)
\]

Pela otimização do lagrangiano associado, temos{[}\^{}1{]} \[
\begin{align*}
\frac{\partial_{1}\tilde{u}}{p_{1}} &= \frac{\partial_{2}\tilde{u}}{p_{2}} \\
\implies \frac{\alpha_{1} x_{1}^{\rho-1}}{p_{1}} &= \frac{\alpha_{2} x_{2}^{\rho -1}}{p_{2}} \\
\therefore x_{2} &= x_{1} \left( \frac{\alpha_{1}}{\alpha_{2}} \frac{p_{2}}{p_{1}} \right)^{\frac{1}{\rho-1}}
\end{align*}
\] Assumindo \(\alpha_{1} = \alpha_{2}\){[}\^{}2{]}, e ressubstituindo
na Restrição Orçamentária, temos \[
\begin{align*}
w &= x_{1}\left( p_{1} + \frac{p_{2}^{1 + 1/(\rho-1)}}{p_{1}^{1/(\rho-1)}} \right) \\
&= x_{1} \frac{(p_{1}^{\rho/(\rho-1)} + p_{2}^{\rho/(\rho-1)})}{p_{1}^{1/(\rho-1)}}
\end{align*}
\]

Defina-se \(\delta \equiv \frac{\rho}{\rho-1}\); portanto,
\(\delta-1 = \frac{1}{\rho-1}\). Estamos analisando
\(\rho \in (-\infty, 1) \iff \delta \in (-\infty, 1)\), onde
\((\rho=1) \leftrightarrow (\delta \to -\infty)\) e vice-versa.

Logo, as demandas walrasianas para a função CES (com pesos iguais) são
\[
\begin{cases}
x_{1} = \frac{p_{1}^{\delta-1}}{p_{1}^\delta + p_{2}^\delta} w \\
x_{2} = \frac{p_{2}^{\delta-1}}{p_{1}^\delta + p_{2}^\delta} w
\end{cases}
\]

Note-se que é possível reescrevê-las como \[
\begin{cases}
x_{1} = \frac{1}{1 + \left( \frac{p_{2}}{p_{1}} \right)^\delta} \frac{w}{p_{1}} \\
x_{2} = \frac{1}{\left( \frac{p_{1}}{p_{2}} \right)^\delta + 1} \frac{w}{p_{2}}
\end{cases}
\]

O caso linear \((\rho=1) \leftrightarrow (\delta \to -\infty)\) depende
dos preços relativos, ficando{[}\^{}3{]} \[
x(p, w) = 
\begin{cases}
\left( \frac{w}{p_{1}}, 0 \right) \,\, &\text{se } p_{1} < p_{2}  \\
\forall \lambda \in [0,1]: \left( \lambda \frac{w}{p}, (1-\lambda) \frac{w}{p} \right) \,\, &\text{se } p_{1} = p_{2} \equiv p \\
\left( 0, \frac{w}{p_{2}} \right) \,\, &\text{se } p_{1} > p_{2}
\end{cases}
\]

A Função de Leontief vem quando
\((\rho \to -\infty) \leftrightarrow (\delta = 1)\): \[
x_{1} = \frac{w}{p_{1} + p_{2}} = x_{2}
\]

\hypertarget{elasticidades-de-substituiuxe7uxe3o}{%
\subsubsection{Elasticidades de
substituição}\label{elasticidades-de-substituiuxe7uxe3o}}

!500

A elasticidade de substituição de \(x_{1}\) por \(x_{2}\) é definida
como: \[
\xi_{12}(p, w) = \frac{ \partial x_{1}/x_{2} }{ \partial p_{1}/p_{2} } \frac{ p_{1}/p_{2}}{x_{1}/x_{2}} 
\] Em verdade, se está falando da elasticidade no tocante à Taxa
Marginal de Substituição: \[
\xi_{12}(p, w) = \frac{ \partial \ln (\frac{x_{1}}{x_{2}}) }{ \partial \ln(TMS_{12}) } 
\]

Note-se que ela geralmente é \textbf{negativa}, pois, conforme se
aumenta o preço relativo \(\frac{p_{1}}{p_{2}}\) (p.~ex. \(p_{1}\)
aumentando conforme \(p_{2}=const\)), a demanda relativa
\(\frac{x_{1}}{x_{2}}\) aumenta{[}\^{}5{]}; portanto, aumentar um
\textbf{diminui} o outro, e a elasticidade fica negativa.

No caso geral da função CES (com \(\alpha_{1} = \alpha_{2}\)), tem-se
que \[
\frac{x_{1}}{x_{2}} = \left( \frac{p_{1}}{p_{2}} \right)^{\delta-1}
\] Derivando com relação a \(\frac{p_{1}}{p_{2}}\), tem-se \[
\frac{ \partial \left( \frac{x_{1}}{x_{2}} \right) }{ \partial \left( \frac{p_{1}}{p_{2}} \right) } = (\delta-1) \left( \frac{p_{1}}{p_{2}} \right)^{\delta-2} 
\] Portanto, a elasticidade de substituição fica \[
\xi_{12}(p,w) = (\delta - 1) \cancel{ \frac{\left( \frac{p_{1}}{p_{2}} \right)^{\delta - 2} \frac{p_{1}}{p_{2}}}{\left( \frac{p_{2}}{p_{1}} \right)^{\delta - 1}} }
\] Reabrindo com relação ao parâmetro original da CES, tem-se \[
\xi_{12}(p,w) = \delta - 1 = \frac{1}{\rho - 1} \left( = - \frac{1}{1 - \rho} \right)
\]

Para a utilidade linear, \(\xi_{12} \to - \infty\); para Cobb-Douglas,
\(\xi_{12} \to -1\); e para Leontief, \(\xi_{12} \to 0\).

\printbibliography

\end{document}
